\subsection{No-Signaling Principle}\label{sec:no-signaling-principle}

Quantum entanglement creates strong correlations between distant particles (i.e., the measurement of one particle instantaneously affects the state of the other, regardless of the distance between them, \autopageref{deepening:physical-meaning-of-entanglement}). However, these correlations cannot be used to transmit information instantaneously.

\highspace
This idea is known as the \textbf{no-signaling principle}: although measurements on \textbf{entangled systems are correlated}, \textbf{they cannot be exploited for faster-than-light communication}.

\highspace
\begin{flushleft}
    \textcolor{Red2}{\faIcon{exclamation-triangle} \textbf{The Communication Problem}}
\end{flushleft}
Suppose Alice and Bob are at opposite ends of the universe. They each possess one qubit of the Bell pair (\autopageref{eq:bell-states}):
\begin{equation*}
    \ket{\Phi^{+}}
    =
    \frac{\ket{00}+\ket{11}}{\sqrt{2}}
\end{equation*}
The first qubit belongs to Alice, while the second belongs to Bob. More precisely, Alice owns the \textbf{first subsystem} of the bipartite state and Bob owns the \textbf{second subsystem}.
\begin{equation*}
    \ket{00} = \ket{0}_{A} \otimes \ket{0}_{B} \qquad \ket{11} = \ket{1}_{A} \otimes \ket{1}_{B}
\end{equation*}
\begin{equation*}
    \ket{\Phi^{+}} = \frac{\ket{0}_{A} \otimes \ket{0}_{B} + \ket{1}_{A} \otimes \ket{1}_{B}}{\sqrt{2}}
\end{equation*}
Because the two qubits are \textbf{entangled}, measuring one qubit instantaneously determines the state of the other. For example:
\begin{itemize}
    \item If Alice measures $\ket{0}$, Bob's qubit immediately collapses to $\ket{0}$;
    \item If Alice measures $\ket{1}$, Bob's qubit immediately collapses to $\ket{1}$.
\end{itemize}
\textcolor{Red2}{\faIcon{exclamation-triangle}} At first sight, this \textbf{seems to violate relativity}. One may wonder:
\begin{quote}
    ``\emph{Can Bob detect instantly whether Alice has already measured her qubit?}''
\end{quote}
If the answer were yes, then \hl{Alice could communicate information to Bob instantaneously}. For example, Alice could decide:
\begin{itemize}
    \item If Alice measures her qubit $\Rightarrow$ transmit bit $1$;
    \item If Alice does not measure her qubit $\Rightarrow$ transmit bit $0$.
\end{itemize}
This \textbf{would allow faster-than-light communication}, contradicting special relativity.

\highspace
The \textbf{no-signaling principle states that this is impossible}: although entanglement creates instantaneous correlations, \hl{Bob cannot locally determine whether Alice has performed a measurement or not.}

\begin{definitionbox}[: No-Signaling Principle]
    The \definition{No-Signaling Principle} states that quantum \textbf{entanglement cannot be used to transmit information instantaneously} between distant observers.

    \highspace
    More precisely, local measurements performed on one subsystem of an entangled quantum system do not allow another observer to determine whether a measurement has already been performed on the distant subsystem.

    \highspace
    Equivalently, although entangled measurements produce strong correlations, the local measurement statistics observed by one party are independent of the actions performed by the other party.
\end{definitionbox}

\begin{proof}
    We want to understand whether Bob can locally detect whether Alice has already measured her qubit or not.

    \highspace
    Alice and Bob share the Bell pair:
    \begin{equation*}
        \ket{\Phi^+}
        =
        \frac{\ket{00}+\ket{11}}{\sqrt{2}}
        =
        \frac{\ket{0}_A \otimes \ket{0}_B+\ket{1}_A \otimes \ket{1}_B}{\sqrt{2}}
    \end{equation*}
    The question is \hl{whether Alice's measurement changes the probabilities observed by Bob}. If Bob's probabilities changed, then \textbf{Alice could use her measurement choice to send information instantaneously}.

    \highspace
    \textbf{Case 1: Alice does not measure.}

    \noindent
    If Alice does not measure her qubit, the global state remains:
    \begin{equation*}
        \ket{\Phi^+}
        =
        \frac{\ket{00}+\ket{11}}{\sqrt{2}}
    \end{equation*}
    If Bob measures his qubit in the computational basis, he obtains:
    \begin{equation*}
        P_B(0)=\frac{1}{2},
        \qquad
        P_B(1)=\frac{1}{2}
    \end{equation*}
    So Bob sees a \textbf{completely random result}.

    \highspace
    \textbf{Case 2: Alice measures.}

    \noindent
    Now suppose Alice measures her qubit first in the computational basis. There are \textbf{two possible outcomes}:
    \begin{itemize}
        \item With probability $\dfrac{1}{2}$, Alice obtains $\ket{0}$, and the global state collapses to:
        \begin{equation*}
            \ket{00}
        \end{equation*}
        In this case, Bob's qubit becomes $\ket{0}$.

        \item With probability $\dfrac{1}{2}$, Alice obtains $\ket{1}$, and the global state collapses to:
        \begin{equation*}
            \ket{11}
        \end{equation*}
        In this case, Bob's qubit becomes $\ket{1}$.
    \end{itemize}
    Therefore, from Bob's point of view:
    \begin{equation*}
        P_B(0)=\frac{1}{2},
        \qquad
        P_B(1)=\frac{1}{2}
    \end{equation*}
    So \textbf{Bob again sees a completely random result}.

    \highspace
    \textbf{Conclusion.}

    \noindent
    In both cases:
    \begin{equation*}
        P_B(0)=P_B(1)=\frac{1}{2}
    \end{equation*}
    Therefore, \hl{Bob's local measurement statistics are the same whether Alice has measured or not}. This means Bob cannot infer Alice's action by measuring only his own qubit. Thus, no usable information is transmitted instantaneously, and the no-signaling principle is preserved.
\end{proof}

\highspace
\textcolor{Green3}{\faIcon{question-circle} \textbf{What are we trying to prove?}} We start with the Bell state:
\begin{equation*}
    \ket{\Phi^+}
    =
    \frac{\ket{00}+\ket{11}}{\sqrt{2}}
    =
    \frac{\ket{0}_A \otimes \ket{0}_B+\ket{1}_A \otimes \ket{1}_B}{\sqrt{2}}
\end{equation*}
Alice has the first qubit, while Bob has the second qubit. Now suppose Alice is billions of light-years away from Bob. The concern is that if Alice measures her qubit, the joint entangled state changes, and Bob's qubit can subsequently be described by a definite state correlated with Alice's outcome. Can Bob detect this change? If yes, then Alice could use her measurement choice to send information to Bob faster than light, for example:
\begin{itemize}
    \item If Alice measures her qubit $\Rightarrow$ transmit bit $1$;
    \item If Alice does not measure her qubit $\Rightarrow$ transmit bit $0$.
\end{itemize}
This would violate the relativity principle.

\highspace
\textcolor{Green3}{\faIcon{question-circle} \textbf{What is the relativity principle?}} It is not a topic of this course, but it is fundamental to understand the motivation behind the no-signaling principle. The \definition{Einstein's Relativity Principle} core idea is: \hl{no information, influence, or signal can travel faster than the speed of light.} This is not merely a technological limitation, but a fundamental law of spacetime.

\highspace
Now, related to quantum mechanics, why is this principle relevant? Imagine Alice is on Earth and Bob is on Mars. Suppose Alice could send a message instantly. Then there would exist observes for whom: Bob receives the message before Alice sends it, which creates a paradox. In some reference frames, the \textbf{effect would occur before the cause}, but this contradicts our understanding of causality. Therefore, Einstein looked at this and thought\footnote{
    We will discuss the EPR paradox later, but it is interesting to note that Einstein historically disliked the idea that measuring one particle could apparently affect another distant particle instantaneously. He famously referred to this phenomenon as ``spukhafte Fernwirkung'' (spooky action at a distance), expressing his discomfort with the idea that quantum mechanics could allow for such non-local effects. As we will see later, Schrödinger coined the term ``entanglement'' to describe this phenomenon, and it has since become a fundamental aspect of quantum mechanics. Nevertheless, Einstein's concerns about entanglement's implications for locality and causality remain a topic of discussion and debate in the foundations of quantum mechanics.
}: ``\emph{wait, did something travel from Alice to Bob faster than light?}''. This is the motivation behind the no-signaling principle: yes, the correlations are real, but Alice cannot choose the outcome she obtains, so she cannot force Bob's qubit to collapse to a specific state. Therefore, Bob cannot detect whether Alice has measured her qubit or not, and no information is transmitted faster than light. In simple terms, \textbf{Entanglement creates correlations, but it does not allow for Communication}. Entanglement makes distant measurement outcomes correlated, but the no-signaling principle guarantees that these correlations cannot be controlled to transmit information.